\scanpage{107}
correspondant
\[
\begin{aligned}
B_D=D^\perp\bigl(&=\text{normalisateur de }D\text{ dans}\\
&\text{l'alg. de Lie }A\text{ i.e.}\\
&\text{ens des }\unsure{u|x}\in A\text{ tels que}\\
&\deleted{[u,x]\in D\ \text{pour }\unsure{\forall|\text{tt}}\,x\in D}\\
&\operatorname{ad}_{\text{Lie}}(D)\subset D\bigr)
\end{aligned}
\]
\[
(B_D^\times=\text{normalisateur de }D\text{ dans }A^\times)
\]

D'où équivalence de (i) avec (iii$'$). D'autre part, (iii$'$) $\Rightarrow$ (iii) trivialemt, il reste à voir que (iii) $\Rightarrow$ (i) -- finalement je ne suis pas si sûr que c'est vrai, sauf sur un corps (ou un anneau réduit) le cas-test décisif est celui où $k=k_0[t]/t^{\unsure{2}}$ \ldots{} (histoire de cohomologie d'une algèbre de Borel \ldots)

\deleted{\uline{Corollaire}.} Notons que $\delta_B=\operatorname{disc}^{\prime}_{\det\mid B}$ est une section de
\[
\delta_B\in(\det B)^{\otimes-2}\left[\underset{\text{iso can}}{\simeq} A/B^{\otimes2}\right]
\]
compte tenu que l'on a un iso can
\[
\det(A)\simeq k
\]
et
\[
\deleted{\unsure{\det A\simeq}}\det B\otimes A/B\simeq\det A\ (\simeq k)
\]
Cet iso dépend d'une convention de signe,
% archive sheet number: 106
